% The paper contains no EPS graphics; disable automatic EPS conversion so the
% documented no-shell-escape build remains warning-free.
\newcommand{\DoNotLoadEpstopdf}{}
\documentclass[11pt]{article}

\usepackage[margin=1in]{geometry}
\usepackage{amsmath,amssymb,amsfonts,amsthm}
\usepackage{array}
\usepackage{bm}
\usepackage{graphicx}
\usepackage{multirow}
\usepackage{multicol}
\usepackage{nicefrac}
\usepackage{paralist}
\usepackage{enumitem}
\usepackage{verbatim}
\usepackage{thm-restate}
\usepackage[normalem]{ulem}
\usepackage[compact]{titlesec}
\usepackage{tikz}
\usetikzlibrary{arrows}
\usetikzlibrary{arrows.meta}
\usetikzlibrary{backgrounds}
\usetikzlibrary{calc}
\usetikzlibrary{decorations.markings}
\usetikzlibrary{fit}
\usetikzlibrary{patterns}
\usetikzlibrary{positioning}
\usetikzlibrary{shapes}

\usepackage{fancybox}
\newenvironment{fminipage}%
  {\begin{Sbox}\begin{minipage}}%
  {\end{minipage}\end{Sbox}\fbox{\TheSbox}}

\newenvironment{algbox}[0]{%
\vskip 0.2in
\noindent
\begin{fminipage}{6.3in}
}{%
\end{fminipage}
\vskip 0.2in
}

\usepackage{tcolorbox}
\tcbuselibrary{skins,breakable}
\tcbset{enhanced jigsaw}
\usepackage{mdframed}
\usepackage{xcolor}
\usepackage{cite}
\usepackage{url}
\usepackage{xspace}
\usepackage[mathscr]{euscript}
\usepackage{nameref}
\usepackage[
  linktocpage=true,
  pagebackref=false,
  colorlinks=true,
  linkcolor=blue,
  citecolor=blue,
  urlcolor=blue,
  bookmarks,
  bookmarksopen,
  bookmarksnumbered
]{hyperref}
\usepackage[noabbrev,nameinlink]{cleveref}

% Explicitly record the substitutions already selected by NFSS for unavailable
% Computer Modern small-cap combinations, avoiding font-substitution warnings.
\DeclareFontShape{OT1}{cmr}{m}{scit}{<->ssub*cmr/m/sc}{}
\DeclareFontShape{OT1}{cmr}{bx}{sc}{<->ssub*cmr/bx/n}{}

\newtheorem{theorem}{Theorem}[section]
\newtheorem{example}{Example}[section]
\newtheorem{lemma}{Lemma}[section]
\newtheorem{proposition}[lemma]{Proposition}
\newtheorem{corollary}[lemma]{Corollary}
\newtheorem{claim}[lemma]{Claim}
\newtheorem{fact}[lemma]{Fact}
\newtheorem{conjecture}[lemma]{Conjecture}
\newtheorem{definition}[lemma]{Definition}
\newtheorem{problem}{Problem}
\newtheorem{remark}[lemma]{Remark}
\newtheorem{observation}[lemma]{Observation}

% Include the parent section in PDF anchor names for counters reset by section.
\renewcommand*{\theHtheorem}{\thesection.\arabic{theorem}}
\renewcommand*{\theHexample}{\thesection.\arabic{example}}
\renewcommand*{\theHlemma}{\thesection.\arabic{lemma}}

\newtheorem*{theorem*}{Theorem}
\newtheorem*{claim*}{Claim}
\newtheorem*{proposition*}{Proposition}
\newtheorem*{lemma*}{Lemma}
\newtheorem*{problem*}{Problem}

\allowdisplaybreaks
\setlength{\parskip}{3pt}

\title{The Instance-Optimal Rate for Bounded Summation under Approximate Differential Privacy}
\author{}
\date{September 1, 2026}

\begin{document}

\pagenumbering{roman}

\maketitle

\begin{abstract}
In the central replacement model, let
$D=(x_1,\ldots,x_n)\in\{0,\ldots,U\}^n$,
$S(D)=\sum_i x_i$, and $m(D)=\max\{1,x_1,\ldots,x_n\}$.  Let
$\rho^*_{\varepsilon,\delta,\beta}(n,U)$ be the infimum, over real-valued
$(\varepsilon,\delta)$-differentially private mechanisms, of the worst-case
ratio $\varepsilon\operatorname{err}_{\beta}(M,D)/m(D)$, where
$\operatorname{err}_{\beta}$ is pointwise $(1-\beta)$-quantile absolute
error.  There are universal constants $C,C_0>0$ such that, for integers
$U\ge2$ and $n\ge1$, $0<\varepsilon\le1$, $0<\beta\le1/3$,
$0\le\delta\le n^{-2}$, and
\[
 n\ge \frac{C_0}{\varepsilon}
 \left(1+\log\log U+\log\frac1\beta\right),
\]
we prove
\[
 \frac1{200}\min\{A_{U,\beta},B_{\varepsilon,\delta}\}
 \le \rho^*_{\varepsilon,\delta,\beta}(n,U)
 \le C\min\{A_{U,\beta},B_{\varepsilon,\delta}\},
\]
where $L_U=1+\lceil\log_2U\rceil$,
$A_{U,\beta}=1+\log(L_U/\beta)$,
$B_{\varepsilon,\delta}=1+\log(1+(e^\varepsilon-1)/\delta)$ for
$\delta>0$, and $B_{\varepsilon,0}=+\infty$.  The upper bound is attained by
a central mechanism outputting a real estimate of $S(D)$.  In the ideal
exact-real and exact-sampling model, its pure-private branch runs in
$O(n\log n+L_U)$ time; its approximate-private branch has deterministic
error $O(m(D)B_{\varepsilon,\delta}/\varepsilon)$ and runs in time polynomial
in $n$, $\log U$, $\log(1/\delta)$, and $\log(1/\varepsilon)$.  The lower
bound has no computational restriction.
\end{abstract}

\noindent\textbf{Note.} This paper was fully AI generated through a harness created by Richard Peng that discovers open problems and runs them through an automated research pipeline.

\clearpage
\tableofcontents
\clearpage

\pagenumbering{arabic}

\section{Introduction}
\label{sec:introduction}

Summation is among the most basic tasks in private data analysis.  Consider
$n$ nonnegative contributions $x_1,\ldots,x_n\in\{0,\ldots,U\}$ and the
query $S(D)=\sum_i x_i$.  Under replacement adjacency, its global
sensitivity is $U$, so the classical Laplace mechanism adds noise at scale
$U/\varepsilon$~\cite{DworkMcSherryNissimSmith2006Calibrating}.  This
worst-case calibration treats a database whose largest entry is $1$ in the
same way as one whose largest entry is $U$.  When $U$ is only a conservative
public bound, an error guarantee governed by the largest value actually
present in the database is substantially more informative.

Clipping makes the issue concrete.  For a threshold $t$, the clipped query
\[
 Q_t(D)=\sum_{i=1}^n\min\{x_i,t\}
\]
has replacement sensitivity $t$, but it incurs the deterministic bias
\[
 S(D)-Q_t(D)=\sum_{i=1}^n(x_i-t)_+.
\]
Thus a low threshold reduces privacy noise at the cost of discarding more of
the sum, while a high threshold has the opposite effect.  This bias--noise
tradeoff also appears in contribution-bounding work, where the best cap can
be characterized with the data held fixed
\cite{AminKuleszaMunozMedinaVassilvitskii2019Contributions}.  In the present
problem, however, the threshold must be chosen from the private data and one
mechanism must give a simultaneous guarantee over every database.

We determine this uniform adaptation cost, up to universal constants, under
central approximate differential privacy.  The formal criterion is as follows.

Let $U\ge2$ and $n\ge1$ be integers, and let databases
$D=(x_1,\ldots,x_n)$ belong to
$\{0,\ldots,U\}^n$, with two databases adjacent when they differ in one
coordinate.  For a real-valued central mechanism $M$, define
\[
 S(D)=\sum_{i=1}^n x_i,
 \qquad
 m(D)=\max\{1,x_1,\ldots,x_n\},
\]
\[
 \operatorname{err}_{\beta}(M,D)
 =\inf\left\{a\ge0:
   \Pr\bigl[|M(D)-S(D)|\le a\bigr]\ge1-\beta\right\},
\]
and
\[
 \rho^*_{\varepsilon,\delta,\beta}(n,U)
 =\inf_M\sup_{D\in\{0,\ldots,U\}^n}
 \frac{\varepsilon\operatorname{err}_{\beta}(M,D)}{m(D)},
\]
where the infimum ranges over all $(\varepsilon,\delta)$-differentially
private mechanisms.  The floor in $m(D)$ keeps the normalization
nondegenerate on the all-zero database.  Unlike a pointwise oracle benchmark,
this definition requires the same private mechanism to attain its guarantee
on all instances.

\subsection{Our result}

Put
\[
 L_U=1+\lceil\log_2U\rceil,
 \qquad
 A_{U,\beta}=1+\log\frac{L_U}{\beta},
\]
and, for $\delta>0$,
\[
 B_{\varepsilon,\delta}
 =1+\log\left(1+\frac{e^\varepsilon-1}{\delta}\right),
 \qquad B_{\varepsilon,0}=+\infty.
\]

\begin{theorem}[Instance-optimal rate for bounded summation]
\label{thm:main}
There are universal constants $C,C_0>0$ such that the following holds for
all integers $U\ge2$ and $n\ge1$.  If
\[
 0<\varepsilon\le1,
 \qquad 0<\beta\le\frac13,
 \qquad 0\le\delta\le n^{-2},
\]
and
\begin{equation}
 n\ge\frac{C_0}{\varepsilon}
 \left(1+\log\log U+\log\frac1\beta\right),
 \label{eq:sample-hypothesis}
\end{equation}
then
\begin{equation}
 \frac1{200}\min\{A_{U,\beta},B_{\varepsilon,\delta}\}
 \le
 \rho^*_{\varepsilon,\delta,\beta}(n,U)
 \le
 C\min\{A_{U,\beta},B_{\varepsilon,\delta}\}.
 \label{eq:main-rate}
\end{equation}
The upper bound is attained by
\hyperref[alg:branch-clip-sum]{\textsc{BranchClipSum}}, which outputs a real
estimate $\widehat S$ satisfying, for every $D$,
\[
 \Pr\left[
  |\widehat S-S(D)|
  \le \frac{C m(D)}{\varepsilon}
  \min\{A_{U,\beta},B_{\varepsilon,\delta}\}
 \right]\ge1-\beta.
\]
Its pure-private branch uses $O(n\log n+L_U)$ exact-real operations.  Its
approximate-private branch has the displayed error on every execution and
runs in time polynomial in $n$, $\log U$, $\log(1/\delta)$, and
$\log(1/\varepsilon)$.  The lower bound in \Cref{eq:main-rate} applies to
all mechanisms, without a computational restriction.
\end{theorem}

The two terms in \Cref{eq:main-rate} describe different privacy regimes.
For $0<\varepsilon\le1$,
\[
 B_{\varepsilon,\delta}
 =\Theta\left(1+\log\left(1+\frac{\varepsilon}{\delta}\right)\right)
 \qquad(\delta>0).
\]
At $\delta=0$, the second term is absent and the theorem gives the
pure-private rate $\Theta(A_{U,\beta})$.  For $\delta>0$, the approximate
branch is smaller precisely when
\begin{equation}
 \delta>
 \frac{(e^\varepsilon-1)\beta}{L_U-\beta}.
 \label{eq:intro-crossover}
\end{equation}
Crossing this threshold selects the better branch; obtaining a strict
asymptotic-order improvement additionally requires
$B_{\varepsilon,\delta}=o(A_{U,\beta})$.

For example, fix $\varepsilon$ and $\beta$, take
$n=\Theta(\log\log U)$ with a sufficiently large implicit constant to meet
\Cref{eq:sample-hypothesis}, and set $\delta=n^{-2}$.  Then
\[
 A_{U,\beta}=\Theta(\log\log U),
 \qquad
 B_{\varepsilon,\delta}=\Theta(\log\log\log U),
\]
and hence
\[
 \rho^*_{\varepsilon,\delta,\beta}(n,U)
 =\Theta(\log\log\log U).
\]
Thus approximate privacy gives a strict improvement over the pure-private
rate in this regime, while the matching lower bound shows that the normalized
ratio still diverges with $U$.

\subsection{Relation to prior work}

For fixed public $n$, estimating a one-dimensional empirical mean and
estimating its sum differ only by multiplication by $n$.  Huang, Liang, and
Yi~\cite{HuangLiangYi2021Mean} study this empirical statistic with
radius- and diameter-adaptive guarantees under $\rho$-zCDP.  Their privacy
notion and their diameter benchmark differ from the standard
$(\varepsilon,\delta)$ replacement-DP and origin-anchored maximum used here.
For pure DP, Dong and Yi~\cite{DongYi2023Universal} give the direct
maximum/radius-based clipping lineage over integer domains and a finite-domain
hard family yielding the double-logarithmic obstruction at constant failure
probability.  The $\delta=0$ specialization of \Cref{thm:main} has this
double-logarithmic dependence and also records the dependence on $\beta$ in
our stated parameter regime.

The same bounded nonnegative sum and realized-maximum accuracy objective have
also been studied under stronger or otherwise altered privacy models.  Dong,
Luo, Fanti, Shi, and Yi~\cite{DongLuoFantiShiYi2024Clipping} obtain the
$A_{U,\beta}$ accuracy scale with a one-round shuffle protocol; their bound
does not use $\delta$ to improve that scale.  Dong, Sun, and
Yi~\cite{DongSunYi2023BetterComposition} give approximate-private
max-contribution upper bounds for vector sums and relational aggregation;
the one-dimensional case is a summation guarantee, but it uses add/remove
user semantics and retains additional terms rather than giving the scalar
replacement-DP transition in \Cref{eq:main-rate}.  A general
approximate-private wrapper for finite-range monotone functions due to
Linder, Raskhodnikova, Smith, and
Steinke~\cite{LinderRaskhodnikovaSmithSteinke2025BlackBox} also implies a
max-scaled sum bound under its add/remove and down-neighborhood formulation;
this is an inferred specialization with a range overhead, rather than the
present max-normalized minimax statement.

Several nearby lines optimize a different object.  Contribution-bound
selection for histograms competes with the best cap in a prescribed family
under user-level privacy~\cite{LiuSureshZhuKairouzGruteser2023Contribution},
while subset-based private mean estimation uses an expected-loss benchmark
defined through deletions from the observed
dataset~\cite{DickKuleszaSunSuresh2023Subset}.  Fixed-sensitivity noise
mechanisms, including truncated Laplace noise, assume that the sensitivity is
already available as a public parameter
\cite{GengDingGuoKumar2020Tradeoff}.  These works respectively illuminate
clipping, alternative notions of instance optimality, and the eventual noise
release, but they do not privately identify the present instance scale and
match it with a lower bound for the same central replacement-DP criterion.

\subsection{Proof ideas}

The upper bound begins by mapping each record to a dyadic scale and retaining
the largest $K$ scales.  This top-$K$ operation is stable under one
replacement.  Any interior point of the retained scales determines a
threshold $t$ satisfying $t<2m(D)$ and clips fewer than $K$ records; therefore
the clipping bias is at most $Km(D)$.  In the pure-private branch, the
exponential mechanism~\cite{McSherryTalwar2007Mechanism} selects an interior
scale with high probability and a Laplace release completes the estimate.  In
the approximate-private branch, an always-interior selector uses one level of
the prefix reduction from the private interior-point
literature~\cite{BunNissimStemmerVadhan2015Thresholds}, after which bounded
truncated Laplace noise is added.  The latter release is in the
fixed-sensitivity tradition of
Geng, Ding, Guo, and Kumar~\cite{GengDingGuoKumar2020Tradeoff}.  Publicly
choosing the branch with the smaller of $A_{U,\beta}$ and
$B_{\varepsilon,\delta}$ gives
\hyperref[alg:branch-clip-sum]{\textsc{BranchClipSum}}.

For the lower bound, we compare the all-zero database with databases
containing the same number of copies of geometrically separated values.
Accuracy forces their output intervals to be disjoint.  Exact approximate
group privacy, with its additive $\delta$ loss retained, then yields both the
$\log(L_U/\beta)$ obstruction and the
$\log(1+(e^\varepsilon-1)/\delta)$ obstruction from a single packing.  A
separate three-point packing supplies the constant floor needed when these
logarithms are small.  Together the arguments give the uniform constant in
the left side of \Cref{eq:main-rate}.

\paragraph{Organization.}
\Cref{sec:preliminaries} fixes the privacy conventions and develops the
ordered-multiset, selection, and prefix tools used in the proof.
\Cref{sec:overview} gives a detailed overview of the two upper-bound branches
and the lower-bound packing.  The full construction and analysis appear in
\Cref{sec:main-proof}.

\section{Preliminaries}
\label{sec:preliminaries}

\paragraph{Conventions and ordered multisets.}
All unqualified logarithms are natural, and $[n]=\{1,\ldots,\allowbreak n\}$.  Multisets
retain multiplicities, and $\uplus$ denotes multiset union.  Two same-size
multisets differ by one replacement if they can be written as
$C\uplus\{x\}$ and $C\uplus\{x'\}$.  All mechanisms below that take
multiset inputs are permutation invariant.  Resource statements use ideal
exact-real arithmetic and exact sampling from the displayed laws; the
minimax statement is information theoretic.

In calls between displayed algorithms, public parameters fixed by the
enclosing specification are ambient and may be suppressed.  Every parameter
whose value differs from its ambient value is displayed explicitly, while
each algorithm's input line records all data and public parameters it uses.

For a finite multiset $Y$ in a totally ordered set and
$0\le K\le |Y|$, let $\operatorname{Top}_K(Y)$ and
$\operatorname{Bot}_K(Y)$ be its $K$ largest and $K$ smallest members,
respectively, with multiplicity.  An \emph{interior point} of a nonempty
multiset $Y$ is an element of its ordered universe lying in
$[\min Y,\max Y]$.  For a point $z$ in the same ordered universe, define
its rank quality by
\begin{equation}
 q_Y(z)=\min\bigl\{
  \#\{y\in Y:y\le z\},\#\{y\in Y:y\ge z\}
 \bigr\}.
 \label{eq:rank-quality}
\end{equation}
One replacement changes every value of $q_Y$ by at most one.  A median of
$Y$ has quality at least $|Y|/2$, whereas a point strictly outside
$[\min Y,\max Y]$ has quality zero.  For every integer $r\ge1$,
\begin{equation}
 q_Y(z)\ge r
 \quad\Longleftrightarrow\quad
 \#\{y\in Y:y\le z\}\ge r
 \ \text{ and }\
 \#\{y\in Y:y\ge z\}\ge r.
 \label{eq:rank-invariant}
\end{equation}
We also write $c_z(Y)=\#\{y\in Y:y\ge z\}$.  For $r\ge1$, a point $z$ is
an \emph{$r$-robust interior point} of $Y$ if
$\min Y\le z$ and $c_z(Y)\ge r$; these conditions imply that $z$ is
interior.

\begin{lemma}[Stability of order truncation]
\label{lem:order-truncation-stability}
If same-size multisets $Y,Y'$ differ by one replacement, then, for every
$0\le K\le|Y|$, the multisets
$\operatorname{Top}_K(Y),\operatorname{Top}_K(Y')$ differ by at most one
replacement, as do
$\operatorname{Bot}_K(Y),\operatorname{Bot}_K(Y')$.  Consequently, applying a
permutation-invariant private mechanism after either truncation preserves
its privacy parameters with respect to the original input.
\end{lemma}

\begin{proof}
Write $Y=C\uplus\{y\}$ and $Y'=C\uplus\{y'\}$.  For
$1\le K<|Y|$, both top-$K$ multisets contain
$\operatorname{Top}_{K-1}(C)$ and one further member; the same statement
holds for the bottom-$K$ multisets.  The cases $K=0$ and $K=|Y|$ are
immediate.
\end{proof}

\paragraph{Privacy tools.}
For $D,D'\in\{0,\ldots,U\}^n$, write $D\sim D'$ when they differ in one
coordinate; for multiset inputs, $D\sim D'$ denotes the one-replacement
relation above.  With either declared adjacency relation, a mechanism $M$ is
$(\varepsilon,\delta)$-differentially private if, for every adjacent pair
and every measurable set $T$,
\begin{equation}
 \Pr[M(D)\in T]\le e^\varepsilon\Pr[M(D')\in T]+\delta.
 \label{eq:dp-definition}
\end{equation}
The replacement $\ell_1$-sensitivity of $f$ taking values in $\mathbb R^k$
is $\sup_{D\sim D'}\|f(D)-f(D')\|_1$; for a scalar query this is its
\emph{replacement sensitivity}.  We use
postprocessing and adaptive basic composition: if, conditional on every
earlier transcript, invocation $j$ is
$(\varepsilon_j,\delta_j)$-private, then the joint transcript is
$(\sum_j\varepsilon_j,\sum_j\delta_j)$-private.  We also use the exact
group-privacy consequence of \Cref{eq:dp-definition}: if $D,D'$ are joined
by a chain of $k$ replacements, then
\begin{equation}
 \Pr[M(D)\in T]
 \le e^{k\varepsilon}\Pr[M(D')\in T]
   +\delta\sum_{j=0}^{k-1}e^{j\varepsilon}.
 \label{eq:group-privacy}
\end{equation}

For $s>0$, $\operatorname{Lap}(s)$ denotes the law with density
$(2s)^{-1}e^{-|w|/s}$.  If $W\sim\operatorname{Lap}(s)$ and $u\ge0$, then
\begin{equation}
 \Pr[|W|>u]=e^{-u/s},
 \qquad
 \Pr[W>u]=\Pr[W<-u]=\tfrac12e^{-u/s}.
 \label{eq:laplace-tail}
\end{equation}
Thus adding independent $\operatorname{Lap}(\Delta/\alpha)$ noise to the
coordinates of a query with $\ell_1$-sensitivity at most $\Delta$ is
$(\alpha,0)$-private.

Suppose a quality function $q_D:\mathcal R\to\mathbb R$ on a finite range
satisfies
$|q_D(z)-q_{D'}(z)|\le\Delta$ for every $D\sim D'$ and
$z\in\mathcal R$, where $\Delta>0$.  The exponential mechanism with
parameter $\alpha$ samples $Z$ according to
\[
 \Pr[Z=z]\ \propto\
 \exp\left(\frac{\alpha q_D(z)}{2\Delta}\right).
\]
It is $(\alpha,0)$-private.  Moreover, for any $z_0\in\mathcal R$ and
$u\ge0$,
\begin{equation}
 \Pr[q_D(Z)\le q_D(z_0)-u]
 \le |\mathcal R|\exp\left(-\frac{\alpha u}{2\Delta}\right).
 \label{eq:exponential-tail}
\end{equation}

For probability laws $P,Q$, their total-variation distance is
$d_{\mathrm{TV}}(P,Q)=\sup_T|P(T)-Q(T)|$; any coupling of random variables
$X\sim P$ and $Y\sim Q$ gives
$d_{\mathrm{TV}}(P,Q)\le\Pr[X\ne Y]$.  If a family $P_D$ is
$(\alpha,\zeta)$-private and
$d_{\mathrm{TV}}(P_D,Q_D)\le\tau$ for every $D$, then, for adjacent
$D,D'$ and every event $T$,
\begin{equation}
 Q_D(T)\le e^\alpha Q_{D'}(T)
       +\zeta+(1+e^\alpha)\tau.
 \label{eq:tv-dp}
\end{equation}

For $\alpha,\Delta>0$ and $0<\zeta\le1/2$, define
\[
 R(\alpha,\zeta,\Delta)
 =\frac{\Delta}{\alpha}
   \log\left(1+\frac{e^\alpha-1}{2\zeta}\right).
\]
The corresponding truncated-Laplace law has density proportional to
$e^{-\alpha|w|/\Delta}$ on
$[-R(\alpha,\zeta,\Delta),\allowbreak R(\alpha,\zeta,\Delta)]$ and density zero
outside this interval.

\paragraph{Dyadic scales and binary prefixes.}
For $x\in\{0,\ldots,U\}$, define the dyadic scale and, for $t\ge0$, the
clipped sum by
\begin{align}
 b(x)&=\left\lceil\log_2\max\{1,x\}\right\rceil
       \in\{0,\ldots,L_U-1\},
 \label{eq:dyadic-scale}\\
 Q_t(D)&=\sum_{i=1}^n\min\{x_i,t\}.
 \label{eq:clipped-sum}
\end{align}

\begin{lemma}[Sensitivity of clipped summation]
\label{lem:clipped-sensitivity}
For every fixed $t\ge0$, the replacement sensitivity of $Q_t$ is at most
$t$.
\end{lemma}

\begin{proof}
A replacement changes one summand, and every summand lies in $[0,t]$.
\end{proof}

For the prefix construction, fix $d\ge2$, put
$\mathcal X=\{0,\ldots,d-1\}$ and $\ell=\lceil\log_2d\rceil$, and let
$\operatorname{bin}_{\ell}(x)\in\{0,1\}^{\ell}$ be the binary
representation of $x\in\mathcal X$ padded with leading zeroes.  If
$x,x'\in\mathcal X$, their longest-common-prefix length is
\[
 \lambda(x,x')=
 \max\left\{s\in\{0,\ldots,\ell\}:
  \operatorname{bin}_{\ell}(x)_{1:s}
  =\operatorname{bin}_{\ell}(x')_{1:s}\right\}.
\]
For $0\le s\le\ell$ and $p\in\{0,1\}^s$, let
\[
 C(p)=\{x\in\mathcal X:
       \operatorname{bin}_{\ell}(x)_{1:s}=p\}.
\]
When $C(p)$ is nonempty, write $a(p)=\min C(p)$ and $b(p)=\max C(p)$;
then $C(p)=[a(p),b(p)]\cap\mathcal X$, also when $d$ is not a power of
two.  For a multiset $Y$ over $\mathcal X$ and a fixed prefix length $s$,
write
\[
 c_Y(p)=|Y\cap C(p)|,
 \qquad
 G_s(Y)=\{p\in\{0,1\}^s:c_Y(p)>0\}
\]
for its prefix counts and its set of observed prefixes.

\section{Overview}
\label{sec:overview}

Put $T=\min\{A_{U,\beta},B_{\varepsilon,\delta}\}$.  The upper construction
and two complementary packing arguments give the matching bounds.  The main
algorithm
\hyperref[alg:branch-clip-sum]{\textsc{BranchClipSum}} receives the database
and the public parameters, chooses publicly between a pure-private and an
approximate-private clipping mechanism, and is $(\varepsilon,\delta)$-private
with pointwise $(1-\beta)$-quantile error $O(m(D)T/\varepsilon)$.  The lower
bound applies to every private mechanism, without a computational
restriction.  The sample hypothesis is used through the normalization
$\varepsilon n\ge64A_{U,\beta}$, which closes the packing case requiring more
than $n$ records and, with $\delta\le n^{-2}$, gives the parameter comparison
needed by the approximate branch.

\paragraph{The common clipping reduction.}
Both upper branches first replace $x_i$ by its dyadic scale $b(x_i)$ and
form
\[
 Y=\operatorname{Top}_K\bigl(\{b(x_i):i\in[n]\}\bigr).
\]
This transformation serves two purposes.  Dyadic scaling reduces threshold
selection from the original universe to only $L_U=O(1+\log U)$ ordered
scales, while retaining the largest $K$ scales makes it enough to locate an
interior point rather than estimate the maximum itself.  Moreover,
top-$K$ order truncation changes by at most one replacement when the input
database does, so a private selector run on $Y$ remains private for the
original database.

The deterministic clipping lemma is the central reduction.  If
$z\in[\min Y,\max Y]$ and $t=2^z$, then
\begin{equation}
 t<2m(D),\qquad
 \#\{i:x_i>t\}\le K-1,
 \qquad
 0\le S(D)-Q_t(D)\le K m(D).
 \label{eq:overview-clipping}
\end{equation}
Thus an interior scale simultaneously keeps the sensitivity of $Q_t$ within
a constant factor of the instance scale and limits the clipping bias.  Once
$Q_t(D)$ is privately released as $\widehat S$, the two contributions obey
\[
 |\widehat S-S(D)|
 \le K m(D)+|\widehat S-Q_t(D)|.
\]
If a prescribed value of $K$ exceeds $n$, the corresponding branch instead
returns zero; the elementary bound
$\varepsilon|S(D)|/m(D)\le\varepsilon n<\varepsilon K$ gives the same
order of error.

\paragraph{The pure-private branch.}
\hyperref[alg:pure-clip-sum]{\textsc{PureClipSum}} takes
$K_{\mathrm{pure}}
 =\left\lceil\frac8\varepsilon\log\frac{2L_U}{\beta}\right\rceil
 =O(A_{U,\beta}/\varepsilon)$.
On the top scales it applies the exponential mechanism, with half the
privacy budget, to the balanced rank quality $q_Y$.  A median of $Y$ has
quality at least $K_{\mathrm{pure}}/2$, whereas every scale outside the
sample range has quality zero, so the selected scale is interior with
probability at least $1-\beta/2$.  Using the remaining privacy budget, it
returns $Q_t(D)+Z$ for
$Z\sim\operatorname{Lap}(2t/\varepsilon)$.  Its standard tail bound and
\Cref{eq:overview-clipping} give normalized $\beta$-error
$O(A_{U,\beta})$.  The required sorting, rank-quality computation, and
clipped summation take $O(n\log n+L_U)$ exact-real operations.

\paragraph{The approximate-private interior-point construction.}
To obtain the smaller $B_{\varepsilon,\delta}$ rate, the proof must avoid the
pure branch's $\log(L_U/\beta)$ cost for making the threshold interior and
must replace the unbounded Laplace tail by a deterministic noise bound.  The
main technical challenge is therefore an approximate-private selector on an
ordered universe of size $d$ whose output is interior on every execution and
whose sample cost depends only doubly logarithmically on $d$.

The starting mechanism,
\hyperref[alg:one-level-prefix]{\textsc{OneLevelPrefix}}, takes a multiset
$S\in\{0,\ldots,d-1\}^K$ and privacy and failure parameters
$(\alpha,\zeta,\eta)$.  It is $(\alpha,\zeta)$-private and, except with
probability $\eta$, returns a robust interior point using
\begin{equation}
 K=O\left(\frac1\alpha\left(
  1+\log\log(4d)+\log\frac1\zeta+\log\frac1\eta
    +\log\frac1\alpha\right)\right)
 \label{eq:overview-one-level-samples}
\end{equation}
records.  Its main intermediate objects preserve the rank information needed
to lift a selected prefix length to a valid point of the original universe.
The mechanism stably bottom-truncates $S$, randomly pairs the retained records, and forms
the multiset of their longest-common-prefix lengths.  An exponential
mechanism chooses a prefix length with both lower and upper rank margins.
The random-pairing separation lemma and the upper margin guarantee that some
binary-prefix cylinder at the next level is well populated.  This ensures
that the
\hyperref[alg:sparse-prefix]{\textsc{SparsePrefix}} subroutine privately
returns an observed cylinder rather than $\bot$ with high probability; the
returned cylinder need not be the populous one.  The subroutine perturbs only
observed prefix counts.  Approximate privacy pays for the possible singleton
bin that can appear or disappear under one replacement, so there is no need
to enumerate all prefixes.  On the same rank events, for any returned
observed cylinder the lower margin and the largest records omitted by bottom
truncation make at least one endpoint robust; a noisy count selects a robust
endpoint with high probability.

The base mechanism is only high-probability interior.  The
\hyperref[alg:sure-interior]{\textsc{SureInterior}} conversion runs two
independent base calls with reduced privacy parameters, returns the first
output that actually lies in $[\min S,\max S]$, and returns $\min S$ only if
both calls fail.  This validation-and-fallback step is data dependent, but
its law is within $O(\eta^2)$ in total variation of the base law conditioned
on being interior.  Taking $\eta$ proportional to
$\min\{\alpha,\zeta\}$ leaves enough privacy slack to absorb the correction.
The result is an $(\alpha,\zeta)$-private, always-interior selector with
the sample bound in \Cref{eq:overview-one-level-samples} with the
$\log(1/\eta)$ term absorbed by
$\log(1/\alpha)+\log(1/\zeta)$.  It runs in time polynomial in $K$,
$\log d$, $\log(1/\zeta)$, and $\log(1/\alpha)$.

\paragraph{The approximate-private branch.}
For $\delta>0$, define
\[
 G=1+\log\log(4L_U)+\log\frac1\delta+\log\frac1\varepsilon,
 \qquad
 K_{\mathrm{app}}=\left\lceil\frac{C_1G}{\varepsilon}\right\rceil.
\]
\hyperref[alg:approx-clip-sum]{\textsc{ApproxClipSum}} applies
\textsc{SureInterior} with parameters $(\varepsilon/2,\delta/2)$ to the
top $K_{\mathrm{app}}$ scales.  This makes
\Cref{eq:overview-clipping} deterministic.  It then invokes
\hyperref[alg:truncated-laplace-release]{\textsc{TruncatedLaplaceRelease}}
on $Q_t$.  For a sensitivity-$\Delta$ scalar query this subroutine uses a
Laplace density truncated to radius
$R(\alpha,\zeta,\Delta)
=(\Delta/\alpha)\log(1+(e^\alpha-1)/(2\zeta))$.
The radius is chosen so that neighboring densities have likelihood ratio at
most $e^\alpha$ on their common support and boundary mass at most $\zeta$;
thus the release is approximately private with deterministically bounded
noise.  With
$(\alpha,\zeta,\Delta)=(\varepsilon/2,\delta/2,t)$, it gives
\begin{equation}
 |Z|\le
 \frac{2t}{\varepsilon}
 \log\left(1+\frac{e^{\varepsilon/2}-1}{\delta}\right)
 \le\frac{4m(D)}{\varepsilon}B_{\varepsilon,\delta}.
 \label{eq:overview-bounded-noise}
\end{equation}
The assumptions $\delta\le n^{-2}$ and
$\varepsilon n\ge64A_{U,\beta}$ imply the parameter comparison
$G=O(B_{\varepsilon,\delta})$.  Therefore the clipping bias
$K_{\mathrm{app}}m(D)$ and the bounded noise in
\Cref{eq:overview-bounded-noise} together give deterministic error
$O(m(D)B_{\varepsilon,\delta}/\varepsilon)$.  This branch runs in time
polynomial in $n$, $\log U$, $\log(1/\delta)$, and
$\log(1/\varepsilon)$.  The public wrapper runs the pure branch when
$\delta=0$ or $A_{U,\beta}\le B_{\varepsilon,\delta}$, and otherwise runs
the approximate branch, proving the upper bound with exactly one invocation.

\paragraph{The matching lower bound.}
Fix any $(\varepsilon,\delta)$-private mechanism $M$, let $R(M)$ be its
worst normalized error, and choose $r>R(M)$.  Set
$a=r/\varepsilon$, $k=\lfloor2a\rfloor+1$, and
$h=1+\lfloor\log_4U\rfloor$.  When $k\le n$, compare the all-zero database
with the $h$ databases having $k$ copies of $4^j$ and zeros elsewhere.
Accuracy places their outputs in pairwise disjoint intervals: the factor
four separates consecutive scales, and $k>2a$ separates all of them from
the interval around zero.  Retaining the additive loss in exact approximate
group privacy and summing these disjoint events under the all-zero database
gives
\begin{equation}
 e^{-k\varepsilon}
 \le\frac{\beta/h+\gamma}{1-\beta+\gamma},
 \qquad
 \gamma=
 \begin{cases}
  \delta/(e^\varepsilon-1),&\delta>0,\\
  0,&\delta=0.
 \end{cases}
 \label{eq:overview-star}
\end{equation}
This single inequality yields the two alternatives
$\log(h/\beta)$ and $\log((e^\varepsilon-1)/\delta)$, which are, up to
universal additive constants, $A_{U,\beta}$ and
$B_{\varepsilon,\delta}$.  If $k>n$, the normalized sample condition gives
the desired lower bound directly.

The additive constants in this logarithmic packing matter when $T$ is
small.  A separate argument assumes $R(M)<1/20$, chooses its own
$r\in(R(M),1/20)$ and corresponding $k$, and packs databases with
$0$, $k$, and $2k$ ones to obtain a contradiction.  This proves the uniform
floor $R(M)\ge1/20$.  Combining that floor with
\Cref{eq:overview-star} gives
$R(M)\ge T/200$.  Taking the infimum over all mechanisms proves the stated
information-theoretic lower bound and completes the match with
\textsc{BranchClipSum}.

\section{Proof of the Main Theorem}
\label{sec:main-proof}

Unless a statement explicitly changes the scope, all parameters in this
section satisfy the hypotheses of \Cref{thm:main}.

\subsection{Scope and parameter normalization}

The floor in $m(D)$ makes the normalized error meaningful at the all-zero
database.  The restriction on $\beta$ is also substantive.

\begin{lemma}[Necessity of a bounded failure probability]
\label{lem:beta-counterexample}
If the restriction $\beta\le1/3$ is dropped and
$1-\beta\le1/(nU+1)$, then
$\rho^*_{\varepsilon,\delta,\beta}(n,U)=0$ for every
$\varepsilon,\delta\ge0$.
\end{lemma}

\begin{proof}
Let $M$ ignore its input and output a uniformly random member of
$\{0,1,\ldots,nU\}$.  This mechanism is $(0,0)$-differentially private.
For every $D$, it outputs $S(D)$ with probability
$1/(nU+1)\ge1-\beta$, so
$\operatorname{err}_{\beta}(M,D)=0$ for every database.
\end{proof}

We next put the sample hypothesis into the form used by both the upper and
lower bounds and relate the exact expression $B_{\varepsilon,\delta}$ to its
more familiar scale.

\begin{lemma}[Parameter normalization]
\label{lem:parameter-bookkeeping}
There is a universal choice of $C_0$ such that
\Cref{eq:sample-hypothesis} implies
\begin{equation}
 \varepsilon n\ge64A_{U,\beta}.
 \label{eq:normalized-sample}
\end{equation}
Moreover, for $0<\varepsilon\le1$ and $\delta>0$,
\[
 B_{\varepsilon,\delta}
 =\Theta\left(1+\log\left(1+\frac{\varepsilon}{\delta}\right)\right)
\]
with universal constants.
\end{lemma}

\begin{proof}
For $U\ge2$, $L_U\le5\log U$, and hence
\[
 A_{U,\beta}
 \le1+\log\log U+\log(1/\beta)+\log5.
\]
Because $\beta\le1/3$, the expression in parentheses in
\Cref{eq:sample-hypothesis} is bounded below by a positive universal
constant.  Increasing $C_0$ therefore gives
\Cref{eq:normalized-sample}.  Finally,
\[
 \varepsilon\le e^\varepsilon-1\le(e-1)\varepsilon
 \qquad(0<\varepsilon\le1),
\]
and substitution in the definition of $B_{\varepsilon,\delta}$ proves the
second assertion.
\end{proof}

\subsection{Stable dyadic clipping}

The following deterministic fact controls the clipping bias and the scale
of the subsequent noisy release.

\begin{lemma}[Bias from an interior dyadic scale]
\label{lem:dyadic-clipping}
Let $y_{[1]}\ge\cdots\ge y_{[n]}$ be the sorted values of $b(x_i)$.  If
$1\le K\le n$, $y_{[K]}\le z\le y_{[1]}$, and $t=2^z$, then
\[
 t<2m(D),
 \qquad
 \#\{i:x_i>t\}\le K-1,
\]
and therefore
\begin{equation}
 0\le S(D)-Q_t(D)\le K m(D).
 \label{eq:clipping-bias}
\end{equation}
\end{lemma}

\begin{proof}
Let $M_D=\max_i x_i$.  Since $z\le b(M_D)$,
\[
 t\le2^{b(M_D)}<2\max\{1,M_D\}=2m(D);
\]
when $M_D\le1$, the sharper equality $t=1$ holds.  If $x_i>2^z$, then
$b(x_i)>z\ge y_{[K]}$.  Fewer than $K$ scales strictly exceed
$y_{[K]}$, which proves the count bound.  Hence
\[
 S(D)-Q_t(D)
 =\sum_i\max\{x_i-t,0\}
 \le \#\{i:x_i>t\}\,M_D
 \le K m(D).
\]
\end{proof}

\subsection{The pure-private branch}

Set
\begin{equation}
 K_{\mathrm{pure}}
 =\left\lceil\frac8\varepsilon\log\frac{2L_U}{\beta}\right\rceil.
 \label{eq:pure-k}
\end{equation}

\begin{algbox}
\textbf{\uline{\textsc{PureClipSum}}}: Pure-private clipping and release.
\phantomsection\label{alg:pure-clip-sum}

\uline{Input}: A database $D\in\{0,\ldots,U\}^n$; public integers
$U\ge2$, $n\ge1$ and public parameters $0<\varepsilon\le1$,
$0<\beta<1$.

\uline{Output}: A real estimate of $S(D)$.

\uline{Guarantee}: $(\varepsilon,0)$-privacy and normalized
$\beta$-error $O(A_{U,\beta})$.

\begin{enumerate}[leftmargin=*,nosep]
  \item If $K_{\mathrm{pure}}>n$, return $0$.
  \item Form
  $Y=\operatorname{Top}_{K_{\mathrm{pure}}}
  (\{b(x_i):i\in[n]\})$ and sample $z\in\{0,\ldots,L_U-1\}$ with
  probability proportional to $\exp(\varepsilon q_Y(z)/4)$.
  \item Set $t=2^z$, independently draw
  $Z\sim\operatorname{Lap}(2t/\varepsilon)$, and return $Q_t(D)+Z$.
\end{enumerate}
\end{algbox}

\begin{lemma}[Pure-private clipping guarantee]
\label{lem:pure-upper}
For integers $U\ge2$, $n\ge1$ and parameters
$0<\varepsilon\le1$, $0<\beta<1$,
\hyperref[alg:pure-clip-sum]{\textsc{PureClipSum}} is
$(\varepsilon,0)$-differentially private and, for every $D$,
\[
 \Pr\left[
  |\text{\textsc{PureClipSum}}(D)-S(D)|
  \le\frac{C m(D)}\varepsilon A_{U,\beta}
 \right]\ge1-\beta
\]
for a universal $C$.  It runs in $O(n\log n+L_U)$ exact-real operations.
\end{lemma}

\begin{proof}
If $K_{\mathrm{pure}}>n$, the output is data independent and
\[
 \frac{\varepsilon|S(D)|}{m(D)}
 \le\varepsilon n<\varepsilon K_{\mathrm{pure}}=O(A_{U,\beta}).
\]
Assume $K_{\mathrm{pure}}\le n$.  By
\Cref{lem:order-truncation-stability} and the sensitivity-one quality
\Cref{eq:rank-quality}, the exponential-mechanism step is
$(\varepsilon/2,0)$-private.  Conditional on every fixed preceding output
$t$, \Cref{lem:clipped-sensitivity} shows that $Q_t$ has sensitivity at most
$t$, so adding $\operatorname{Lap}(2t/\varepsilon)$ is also
$(\varepsilon/2,0)$-private.  Adaptive composition proves
$(\varepsilon,0)$-privacy.

Some median of $Y$ has quality at least $K_{\mathrm{pure}}/2$, while every
point strictly outside $[\min Y,\max Y]$ has quality zero.  Thus
\Cref{eq:exponential-tail} gives
\[
 \Pr[z\notin[\min Y,\max Y]]
 \le L_Ue^{-\varepsilon K_{\mathrm{pure}}/8}
 \le\frac\beta2.
\]
The Laplace tail \Cref{eq:laplace-tail} also gives
\[
 \Pr\left[|Z|>\frac{2t}{\varepsilon}\log\frac2\beta\right]
 \le\frac\beta2.
\]
On the complementary events, \Cref{lem:dyadic-clipping} yields
\[
 |Q_t(D)+Z-S(D)|
 \le K_{\mathrm{pure}}m(D)
   +\frac{4m(D)}\varepsilon\log\frac2\beta
 \le\frac{C m(D)}\varepsilon A_{U,\beta}.
\]
A union bound proves accuracy.

For the resource bound, compute and sort the $n$ scales, fill an
$L_U$-entry histogram of the selected scales, and use one prefix and one
suffix pass to compute all qualities.  These operations cost
$O(n\log n+L_U)$.  Sampling the resulting categorical law costs $O(L_U)$
exact-real operations, evaluating $Q_t$ costs $O(n)$, and the Laplace draw
costs $O(1)$.
\end{proof}

\subsection{Bounded noise for approximate privacy}

The approximate-private branch uses the truncated-Laplace law from the
preliminaries.  Writing
$\lambda=\alpha/\Delta$, $R=R(\alpha,\zeta,\Delta)$, and
$E=e^{-\lambda R}$, its magnitude is
\[
 -\lambda^{-1}\log\bigl(1-(1-E)V\bigr)
\]
for $V$ uniform on $[0,1]$, and its sign is an independent fair sign.

\begin{algbox}
\textbf{\uline{\textsc{TruncatedLaplaceRelease}}}: Bounded-noise scalar
release.  \phantomsection\label{alg:truncated-laplace-release}

\uline{Input}: A database $D$, a scalar query $q$, and public parameters
$\Delta>0$, $\alpha>0$, and $0<\zeta\le1/2$, where $q$ has replacement
sensitivity at most $\Delta$.

\uline{Output}: A real release of $q(D)$.

\uline{Guarantee}: $(\alpha,\zeta)$-privacy and deterministic error at most
$R(\alpha,\zeta,\Delta)$.

\begin{enumerate}[leftmargin=*,nosep]
  \item Put $R=R(\alpha,\zeta,\Delta)$ and $\lambda=\alpha/\Delta$.
  Draw $V$ uniformly from $[0,1]$ and an independent fair sign
  $\sigma\in\{-1,1\}$.
  \item Set
  $W=-\sigma\lambda^{-1}
  \log(1-(1-e^{-\lambda R})V)$ and return $q(D)+W$.
\end{enumerate}
\end{algbox}

\begin{lemma}[Privacy and bounded error of truncated Laplace release]
\label{lem:truncated-laplace}
For every input in its specification,
\hyperref[alg:truncated-laplace-release]{\textsc{TruncatedLaplaceRelease}}
is $(\alpha,\zeta)$-differentially private and has absolute error at most
$R(\alpha,\zeta,\Delta)$ on every execution.  If evaluating $q(D)$ takes
time $T_q(D)$, its running time is $T_q(D)+O(1)$ in the exact-real model.
\end{lemma}

\begin{proof}
Consider adjacent query values $q(D)\le q(D')=q(D)+s$ with
$0\le s\le\Delta$.  Because $\zeta\le1/2$,
$R(\alpha,\zeta,\Delta)\ge\Delta$.  On the overlap of the two shifted
supports, the density ratio is at most
$e^{\alpha s/\Delta}\le e^\alpha$.  The part of either support outside the
other has probability at most the mass of a boundary interval of width
$\Delta$.  With $R=R(\alpha,\zeta,\Delta)$, this mass is
\[
 \frac{e^{-\alpha(R-\Delta)/\Delta}-e^{-\alpha R/\Delta}}
      {2(1-e^{-\alpha R/\Delta})}
 =\frac{e^\alpha-1}{2(e^{\alpha R/\Delta}-1)}
 =\zeta.
\]
Splitting an arbitrary event into its overlap and nonoverlap parts proves
the privacy inequality in both directions.  The support gives the
deterministic error bound.  The inverse-CDF representation in the algorithm
uses a constant number of exact-real operations and random draws in addition
to evaluating $q(D)$.
\end{proof}

\subsection{A rank-preserving interior-point construction}

We now construct an always-interior selector on the ordered universe
$\mathcal X=\{0,\ldots,d-1\}$ for $d\ge2$, using the binary-prefix notation
from the preliminaries.  The next selector inspects only prefixes that occur
in its input.  Approximate
privacy pays for the possible singleton prefix that appears or disappears
under a replacement, avoiding enumeration of all $2^s$ prefixes.

For $0<\rho,\zeta,\nu\le1/4$, define
\begin{equation}
 b_\rho=\frac2\rho,
 \qquad
 \tau_{\rho,\zeta}
 =1+\frac2\rho\log\frac{1+e^\rho}{2\zeta},
 \qquad
 h_{\rho,\zeta,\nu}
 =\left\lceil
   \tau_{\rho,\zeta}+\frac2\rho\log\frac1{2\nu}
  \right\rceil.
 \label{eq:sparse-thresholds}
\end{equation}
\begin{algbox}
\textbf{\uline{\textsc{SparsePrefix}}}: Private selection from observed
prefixes.  \phantomsection\label{alg:sparse-prefix}

\uline{Input}: A multiset $D$ over $\mathcal X$, a prefix length
$s\in\{0,\ldots,\ell\}$, and parameters
$0<\rho,\zeta,\nu\le1/4$, in this order.

\uline{Output}: An observed length-$s$ prefix, or $\bot$.

\uline{Guarantee}: $(\rho,\zeta)$-privacy; if some prefix occurs at least
$h_{\rho,\zeta,\nu}$ times, the output is an observed prefix with probability
at least $1-\nu$.

\begin{enumerate}[leftmargin=*,nosep]
  \item Form the observed-prefix histogram
  $\{(p,c_D(p)):p\in G_s(D)\}$.
  \item For each $p\in G_s(D)$, independently draw
  $E_p\sim\operatorname{Lap}(b_\rho)$ and declare $p$ active when
  $c_D(p)+E_p\ge\tau_{\rho,\zeta}$.
  \item Return the lexicographically first active prefix, or return $\bot$
  if none is active.
\end{enumerate}
\end{algbox}

\begin{lemma}[Private sparse-prefix selection]
\label{lem:sparse-prefix}
\hyperref[alg:sparse-prefix]{\textsc{SparsePrefix}} has the privacy and
utility guarantees in its specification.  It runs in $O(|D|\ell)$ bit
operations and uses at most $|D|$ exact Laplace samples after the histogram
is formed.
\end{lemma}

\begin{proof}
Fix adjacent multisets $D,D'$ and let
$H=G_s(D)\cap G_s(D')$.  Each of
$G_s(D)\setminus H$ and $G_s(D')\setminus H$ has cardinality at most one,
and a prefix in either difference has count one in the multiset where it
occurs.  Put $u=\zeta/(1+e^\rho)$.  Such a singleton prefix is activated with
probability
\begin{equation}
 \Pr[1+E_p\ge\tau_{\rho,\zeta}]
 =\frac12\exp\left(-\frac{\tau_{\rho,\zeta}-1}{2/\rho}\right)
 =\frac{\zeta}{1+e^\rho}=u.
 \label{eq:singleton-tail}
\end{equation}

Let $Q_D$ be the same procedure restricted to the common prefix set $H$,
using the counts from $D$, and define $Q_{D'}$ analogously.  Coupling noises
on $H$ shows that the output law of \textsc{SparsePrefix}$(D)$ is within
total variation $u$ of that of $Q_D$, and likewise for $D'$.  On $H$, the
two count vectors have $\ell_1$-distance at most two.  Adding independent
$\operatorname{Lap}(2/\rho)$ noise therefore changes their joint density by
a factor at most $e^\rho$.  Thresholding and selecting the first active
coordinate are postprocessing, so $Q_D$ and $Q_{D'}$ satisfy pure
$\rho$-privacy.  Consequently, for every output event $T$,
\[
 \begin{aligned}
 \Pr[\text{\textsc{SparsePrefix}}(D)\in T]
 &\le \Pr[Q_D\in T]+u\\
 &\le e^\rho\Pr[Q_{D'}\in T]+u\\
 &\le e^\rho\Pr[\text{\textsc{SparsePrefix}}(D')\in T]
       +(1+e^\rho)u,
 \end{aligned}
\]
which is the $(\rho,\zeta)$-privacy inequality.  Interchanging $D,D'$ gives
the reverse direction.

If $c_D(p)\ge h_{\rho,\zeta,\nu}$ for some $p$, then this prefix is inactive
only if
\[
 E_p\le-\frac2\rho\log\frac1{2\nu},
\]
an event of probability at most $\nu$.  Whenever it is active, the algorithm
returns some member of $G_s(D)$.  Finally, a binary-trie scan constructs the
at most $|D|$ observed bins and their counts in $O(|D|\ell)$ bit operations;
the remaining resource claims follow directly from the algorithm.
\end{proof}

To turn a privately selected prefix length into a useful cylinder, we need a
random pairing to contain a pair that is well separated in input rank.

\begin{lemma}[Random-pairing separation]
\label{lem:pairing}
Let $N$ be even, let $(\Pi_1,\ldots,\Pi_N)$ be a uniformly random
permutation of $\{1,\ldots,N\}$, and pair consecutive positions.  For every
integer $r\ge1$,
\begin{equation}
 \Pr\left[
  \#\left\{j:|\Pi_{2j-1}-\Pi_{2j}|\le r/12\right\}\ge r
 \right]\le2^{-r}.
 \label{eq:pairing-bound}
\end{equation}
\end{lemma}

\begin{proof}
If $r>N/2$, the event is empty.  Otherwise call a pair short when its labels
differ by at most $r/12$.  From a permutation with at least $r$ short pairs,
select canonically the first $r$ of them.  There are at most $\binom Nr$
choices for the labels in their first positions, and each such label has at
most $r/6$ possible partner labels.  Compressing the selected directed pairs
into blocks leaves $N-r$ distinct objects, which can be ordered in at most
$(N-r)!$ ways.  This gives an injective encoding after the canonical
selection, and hence the probability is at most
\[
 \frac{\binom Nr(r/6)^r(N-r)!}{N!}
 =\frac{(r/6)^r}{r!}
 \le\left(\frac e6\right)^r<2^{-r},
\]
where $r!\ge(r/e)^r$.
\end{proof}

We now choose public parameters that provide both lower and upper rank
margins.  For $0<\alpha,\zeta,\eta\le1/4$, put
\begin{equation}
 \rho=\frac\alpha3,
 \qquad \nu=\frac\eta4,
 \qquad h=h_{\rho,\zeta,\nu},
 \label{eq:one-level-basic-parameters}
\end{equation}
\begin{equation}
 r=\left\lceil\max\left\{
  24h,
  \frac2\rho\log\frac1{2\nu},
  \frac{\log(1/\nu)}{\log2},
  12
 \right\}\right\rceil,
 \label{eq:rank-guard}
\end{equation}
and
\begin{equation}
 K_0=6r+\frac8\rho\log\frac{\ell+1}{\nu}+2.
 \label{eq:exact-sample-threshold}
\end{equation}
For an input size $K\ge K_0$, define
\begin{equation}
 m=\left\lfloor\frac{K-2r}{2}\right\rfloor,
 \qquad N=2m.
 \label{eq:bottom-size}
\end{equation}
These definitions ensure
\begin{equation}
 K-N\ge2r,
 \qquad
 m\ge2r+\frac4\rho\log\frac{\ell+1}{\nu}.
 \label{eq:pair-count-lower}
\end{equation}
For the prefix-length multiset $W=(w_1,\ldots,w_m)$, we use the rank
quality $q_W$ and the two rank margins in
\Cref{eq:rank-quality,eq:rank-invariant}.

\begin{algbox}
\textbf{\uline{\textsc{OneLevelPrefix}}}: A robust one-level prefix lift.
\phantomsection\label{alg:one-level-prefix}

\uline{Input}: A multiset $S\in\{0,\ldots,d-1\}^K$ and parameters
$\alpha,\zeta,\eta$ satisfying $0<\alpha,\zeta,\eta\le1/4$ and
$K\ge K_0$, in this order.

\uline{Output}: A point of $\{0,\ldots,d-1\}$.

\uline{Guarantee}: The mechanism is permutation invariant and
$(\alpha,\zeta)$-private; except with probability $\eta$, its output is an
$r$-robust interior point of $S$.

\begin{enumerate}[leftmargin=*,nosep]
  \item Form $B=\operatorname{Bot}_N(S)$, uniformly order its labelled
  occurrences as $(y_1,\ldots,y_N)$, and set
  $w_j=\lambda(y_{2j-1},y_{2j})$ for $j\in[m]$.
  \item Sample $z\in\{0,\ldots,\ell\}$ with probability proportional to
  $\exp(\rho q_W(z)/2)$, where $W=(w_1,\ldots,w_m)$.
  \item Set $s=\min\{z+1,\ell\}$ and
  $p=\text{\textsc{SparsePrefix}}(B,s;\rho=\rho,\zeta=\zeta,\nu=\nu)$.
  If $p=\bot$, return $0$;
  otherwise set $a=a(p)$ and $b=b(p)$.
  \item If $z=\ell$, return $a$.  Otherwise compute $c_b(S)$, draw
  $E\sim\operatorname{Lap}(1/\rho)$, and return $b$ when
  $c_b(S)+E\ge3r/2$, or $a$ otherwise.
\end{enumerate}
\end{algbox}

\begin{lemma}[One-level robust prefix lift]
\label{lem:one-level-prefix}
For every input satisfying its specification,
\hyperref[alg:one-level-prefix]{\textsc{OneLevelPrefix}} is permutation
invariant and $(\alpha,\zeta)$-differentially private.  With probability at
least $1-\eta$, its output $x$ is $r$-robust, namely
\[
 \min S\le x,
 \qquad
 \#\{y\in S:y\ge x\}\ge r,
\]
and in particular is an interior point.  Its running time is polynomial in
$K$, $\log d$, $\log(1/\zeta)$, and $\log(1/\eta)$ in the exact-real model.
\end{lemma}

\begin{proof}
We begin with privacy.  By \Cref{lem:order-truncation-stability}, the bottom
truncations of adjacent inputs differ by at most one replacement.  Match
their common labelled occurrences and their possible exceptional
occurrences, and place matched occurrences in the same positions of the two
uniform random orderings.  This couples the orderings so that they differ in
at most one position, and their paired LCP multisets therefore differ in at
most one coordinate.  Since \Cref{eq:rank-quality} has sensitivity one,
the exponential mechanism selecting $z$ is pure $\rho$-private.

Conditional on every selected $z$, the prefix length $s$ is fixed, and
\Cref{lem:order-truncation-stability,lem:sparse-prefix} make the
sparse-prefix call
$(\rho,\zeta)$-private as a mechanism of the original input.  Conditional on
every transcript $(z,p)$, the final count $c_b(S)$ has replacement
sensitivity one; adding $\operatorname{Lap}(1/\rho)$ before thresholding is
pure $\rho$-private.  The remaining branches are postprocessing.  Adaptive
composition gives
\begin{equation}
 (\rho,0)+(\rho,\zeta)+(\rho,0)=(\alpha,\zeta).
 \label{eq:one-level-composition}
\end{equation}

We turn to utility.  A median of $W$ has rank quality at least $m/2$.
The exponential tail \Cref{eq:exponential-tail} and
\Cref{eq:pair-count-lower} imply,
conditional on every $W$,
\begin{equation}
 \Pr[q_W(z)<r\mid W]
 \le(\ell+1)
   \exp\left[-\frac\rho2\left(\frac m2-r\right)\right]
 \le\nu.
 \label{eq:robust-em-utility}
\end{equation}
Thus both rank margins in \Cref{eq:rank-invariant} hold except with
probability $\nu$.

Label the occurrences of $B$ by $1,\ldots,N$ in nondecreasing value order,
breaking ties arbitrarily.  The random ordering induces a uniform
permutation of these ranks.  By \Cref{lem:pairing} and the definition of
$r$, except with probability at most
\begin{equation}
 2^{-r}\le\nu,
 \label{eq:pairing-failure}
\end{equation}
fewer than $r$ pairs have rank distance at most $r/12$.  Suppose this event
and $q_W(z)\ge r$ both hold.  Among the at least $r$ pairs with LCP length at
least $z$, one pair has rank distance greater than $r/12$.  More than
$r/12$ records of $B$ lie between its endpoints in value order, inclusively,
and all share their first $z$ bits.  If $z<\ell$, at least half share one of
the two possible next bits; if $z=\ell$, all have the full common binary
representation.  Hence some prefix of length
$s=\min\{z+1,\ell\}$ occurs more than $r/24\ge h$ times in $B$.

By \Cref{lem:sparse-prefix}, the call
$\text{\textsc{SparsePrefix}}(B,s;\rho=\rho,\zeta=\zeta,\nu=\nu)$ returns an observed prefix $p$
except with probability $\nu$.  Fix a witness
$y^*\in B\cap C(p)$, so
\begin{equation}
 a\le y^*\le b.
 \label{eq:cylinder-witness}
\end{equation}
Every one of the at least $K-N\ge2r$ omitted largest records is at least
every member of $B$.  If $z=\ell$, then $a=b=y^*$, so the returned value has
a lower witness and at least $2r$ records at or above it.

Suppose $z<\ell$.  The lower-rank margin in
\Cref{eq:rank-invariant} supplies a pair whose LCP length is at most $z$.
The two endpoints cannot both lie in the length-$(z+1)$ cylinder $[a,b]$.
Thus either its smaller endpoint is below $a$, or its larger endpoint is
above $b$.  In the former case, $a$ has a lower witness and at least $2r$
inputs at or above it.  In the latter case, $b$ has the lower witness
\Cref{eq:cylinder-witness}, and the omitted records show that at least
$2r$ inputs are at or above $b$.  At least one endpoint is therefore robust.

It remains to analyze the noisy choice.  The point $a$ always has at least
$2r$ inputs at or above it, and it is robust precisely when some input is at
most $a$.  The point $b$ always has the witness $y^*\le b$, and it is robust
whenever $c_b(S)\ge r$.  If $c_b(S)<r$, choosing $b$ requires
$E\ge r/2$ and has probability at most
\begin{equation}
 \Pr[E\ge r/2]\le\tfrac12e^{-\rho r/2}\le\nu.
 \label{eq:noisy-endpoint-failure}
\end{equation}
If $a$ is not robust, the separating pair must instead have an endpoint
above $b$, so $c_b(S)\ge2r$; choosing $a$ then requires $E<-r/2$ and has the
same probability bound.  If both endpoints are robust, either choice works.

The failures in \Cref{eq:robust-em-utility},
\Cref{eq:pairing-failure}, the sparse-prefix call, and
\Cref{eq:noisy-endpoint-failure} each have probability at most $\nu$.
Their union has probability at most $4\nu=\eta$, proving the claimed
invariant.

Finally, sorting costs $O(K\log K)$ comparisons; computing the LCP values
and the sparse observed-prefix histogram costs $O(K\ell)$ bit operations;
histograms compute all $\ell+1$ balanced qualities in $O(K+\ell)$
exact-real operations; and the endpoint count costs $O(K)$.  Since
$\ell=O(\log d)$ and all public thresholds use only the displayed
logarithms, the running time has the asserted polynomial bound.  Uniformly
ordering labelled occurrences also shows permutation invariance.
\end{proof}

The preceding exact sample threshold has the asymptotic form needed by the
summation algorithm.

\begin{lemma}[One-level interior-point sample bound]
\label{lem:ip-sample-bound}
Let
\[
 \Xi=1+\log\log(4d)+\log\frac1\zeta
       +\log\frac1\eta+\log\frac1\alpha.
\]
For $d\ge2$ and $0<\alpha,\zeta,\eta\le1/4$, the parameters defined in
\Cref{eq:one-level-basic-parameters,eq:rank-guard,eq:exact-sample-threshold}
satisfy
\[
 K_0\le\frac{2256}{\alpha}\Xi.
\]
Consequently, $K\ge3000\Xi/\alpha$ implies the precondition of
\hyperref[alg:one-level-prefix]{\textsc{OneLevelPrefix}}.
\end{lemma}

\begin{proof}
Since $\rho=\alpha/3\le1/12$ and $\nu=\eta/4$,
\[
 \log\frac{1+e^\rho}{2\zeta}\le1+\log\frac1\zeta,
 \qquad
 \log\frac1{2\nu}=\log\frac2\eta\le1+\log\frac1\eta.
\]
Thus \Cref{eq:sparse-thresholds} gives $h\le5\Xi/\rho$.  Each of the
other three entries in the maximum in \Cref{eq:rank-guard} is at most
$120\Xi/\rho$, as is $24h$, and therefore
\begin{equation}
 r\le\frac{121\Xi}{\rho}.
 \label{eq:r-comparison}
\end{equation}
Moreover,
\[
 \ell+1\le\log_2(4d),
 \qquad
 \log\frac{\ell+1}{\nu}\le3\Xi.
\]
Substituting these bounds into \Cref{eq:exact-sample-threshold}, and using
$2\le2\Xi/\rho$, yields
\[
 K_0\le\frac{(6\cdot121+24+2)\Xi}{\rho}
 =\frac{752\Xi}{\rho}
 =\frac{2256\Xi}{\alpha}.
\]
\end{proof}

We next convert the high-probability interior point into an output that is
interior on every execution.  For $0<\alpha,\zeta\le1/2$, define
\begin{equation}
 \alpha_0=\frac\alpha4,
 \qquad
 \zeta_0=\frac\zeta{16},
 \qquad
 \eta=\frac{\min\{\alpha,\zeta\}}{64}.
 \label{eq:sure-parameters}
\end{equation}
Let
\begin{equation}
 K_{\mathrm{sure}}
 =\frac{3000}{\alpha_0}\left(
  1+\log\log(4d)+\log\frac1{\zeta_0}
  +\log\frac1\eta+\log\frac1{\alpha_0}
 \right).
 \label{eq:sure-exact-threshold}
\end{equation}
By \Cref{lem:one-level-prefix,lem:ip-sample-bound}, running
\textsc{OneLevelPrefix} with parameters
$(\alpha_0,\zeta_0,\eta)$ on at least $K_{\mathrm{sure}}$ records is
$(\alpha_0,\zeta_0)$-private and returns an interior point with probability
at least $1-\eta$.

\begin{algbox}
\textbf{\uline{\textsc{SureInterior}}}: Conversion to an always-interior
selector.  \phantomsection\label{alg:sure-interior}

\uline{Input}: A multiset $Y\in\{0,\ldots,d-1\}^K$ and parameters
$0<\alpha,\zeta\le1/2$, in this order, with $K\ge K_{\mathrm{sure}}$.

\uline{Output}: A point of $\{0,\ldots,d-1\}$.

\uline{Guarantee}: $(\alpha,\zeta)$-privacy and output in
$[\min Y,\max Y]$ on every execution.

\begin{enumerate}[leftmargin=*,nosep]
  \item Independently compute
  $Z_1=\text{\textsc{OneLevelPrefix}}
  (Y;\alpha=\alpha_0,\zeta=\zeta_0,\eta=\eta)$ and
  $Z_2=\text{\textsc{OneLevelPrefix}}
  (Y;\alpha=\alpha_0,\zeta=\zeta_0,\eta=\eta)$.
  \item Return the first $Z_j$ that belongs to $[\min Y,\max Y]$; if neither
  does, return $\min Y$.
\end{enumerate}
\end{algbox}

\begin{lemma}[Always-interior conversion]
\label{lem:sure-interior}
There is a universal $C$ such that
\hyperref[alg:sure-interior]{\textsc{SureInterior}} is an
$(\alpha,\zeta)$-differentially private, always-interior mechanism whenever
\begin{equation}
 K\ge\frac{C}{\alpha}\left(
  1+\log\log(4d)+\log\frac1\zeta+\log\frac1\alpha
 \right).
 \label{eq:sure-sample}
\end{equation}
If $T_{\mathrm{base}}$ is the running time of one base call with the
parameters in \Cref{eq:sure-parameters}, then the running time is
$2T_{\mathrm{base}}+O(K)$.  It is polynomial in $K$, $\log d$,
$\log(1/\zeta)$, and $\log(1/\alpha)$.
\end{lemma}

\begin{proof}
The parameters in \Cref{eq:sure-parameters} lie in the range required by
\Cref{lem:one-level-prefix}, and
\[
 \log(1/\eta)
 \le\log64+\log(1/\alpha)+\log(1/\zeta).
\]
Thus a universal choice of $C$ in \Cref{eq:sure-sample} makes its
right-hand side at least $K_{\mathrm{sure}}$.

Let $P_Y$ be the output law of one base call, let
$G_Y=[\min Y,\max Y]$, and let $C_Y$ be $P_Y$ conditioned on $G_Y$.  For
every measurable $T$, base privacy and base failure probability at most
$\eta$ imply
\[
 \begin{aligned}
 C_Y(T)
 &\le\frac{P_Y(T)}{1-\eta}\\
 &\le\frac{e^{\alpha_0}P_{Y'}(T)+\zeta_0}{1-\eta}\\
 &\le\frac{e^{\alpha_0}C_{Y'}(T)+e^{\alpha_0}\eta+\zeta_0}
          {1-\eta}.
 \end{aligned}
\]
Hence the family $C_Y$ is $(\alpha_C,\zeta_C)$-private for
\begin{equation}
 \alpha_C=\alpha_0+\log\frac1{1-\eta},
 \qquad
 \zeta_C=\frac{\zeta_0+e^{\alpha_0}\eta}{1-\eta}.
 \label{eq:conditioned-privacy}
\end{equation}

For $x\in\mathcal X$, let $\delta_x$ denote the point-mass law at $x$.
If $p_Y=P_Y(G_Y^c)$, the output law $Q_Y$ of \textsc{SureInterior} is
\[
 Q_Y=(1-p_Y^2)C_Y+p_Y^2\,\delta_{\min Y},
\]
so $d_{\mathrm{TV}}(Q_Y,C_Y)\le\eta^2$.  Apply \Cref{eq:tv-dp} with
$\tau=\eta^2$ and the parameters in
\Cref{eq:conditioned-privacy}.  Since $\eta\le1/128$ and
$\eta\le\alpha/64$, we have
$\log(1/(1-\eta))\le2\eta$ and therefore
\[
 \alpha_C\le\alpha/4+2\eta
 \le\alpha/4+\alpha/32<\alpha.
\]
Also $\eta\le\zeta/64$, $e^{\alpha_0}\le e^{1/8}<8/7$, and therefore
\[
 \zeta_C
 \le\frac{128}{127}\left(\frac1{16}+\frac{8/7}{64}\right)\zeta
 <\frac\zeta{12}.
\]
Finally, $\eta^2\le\zeta/8192$ and $\alpha_C<1/2$, so
\[
 \zeta_C+(1+e^{\alpha_C})\eta^2<\zeta.
\]
This proves $(\alpha,\zeta)$-privacy.  The output is interior by
construction.  One scan computes $\min Y$ and $\max Y$ in $O(K)$ time, and
the algorithm makes exactly two base calls; the resource bounds now follow
from \Cref{lem:one-level-prefix}.
\end{proof}

\subsection{The approximate-private branch}

It remains to compare the sample requirement of the interior-point
construction with the target error rate.  For $\delta>0$, put
\begin{equation}
 G=1+\log\log(4L_U)+\log\frac1\delta+\log\frac1\varepsilon.
 \label{eq:G-definition}
\end{equation}

\begin{lemma}[Approximate-branch parameter comparison]
\label{lem:parameter-comparison}
Under \Cref{eq:sample-hypothesis}, if $0<\delta\le n^{-2}$, then
\begin{equation}
 G\le C B_{\varepsilon,\delta}
 \label{eq:G-versus-B}
\end{equation}
for a universal $C$.
\end{lemma}

\begin{proof}
Put $h=1+\log L_U$.  By \Cref{lem:parameter-bookkeeping},
$\varepsilon n\ge64A_{U,\beta}\ge64h$.  Write
\[
 u=\log\frac\varepsilon\delta,
 \qquad
 e_0=\log\frac1\varepsilon.
\]
Since $\delta\le n^{-2}$,
\[
 \frac\varepsilon\delta
 \ge\varepsilon n^2
 =\frac{(\varepsilon n)^2}{\varepsilon}
 \ge\frac{(64h)^2}{\varepsilon}.
\]
Consequently $u\ge e_0$, and $e^\varepsilon-1\ge\varepsilon$ gives
$B_{\varepsilon,\delta}\ge1+u$.  Furthermore,
\[
 \log\frac1\delta+\log\frac1\varepsilon
 =(u+e_0)+e_0\le3u\le3B_{\varepsilon,\delta}.
\]
Finally, $\log(4L_U)\le2h$, whereas $u\ge2\log(64h)$, and hence
$1+\log\log(4L_U)=O(1+u)$.  Combining the estimates proves the claim.
\end{proof}

Fix a sufficiently large universal $C_1$ and define
\begin{equation}
 K_{\mathrm{app}}=\left\lceil\frac{C_1G}{\varepsilon}\right\rceil.
 \label{eq:approx-k}
\end{equation}
The choice of $C_1$ ensures that \Cref{eq:sure-sample} holds with
$d=L_U$, $\alpha=\varepsilon/2$, and $\zeta=\delta/2$ whenever
$K_{\mathrm{app}}\le n$.

\begin{algbox}
\textbf{\uline{\textsc{ApproxClipSum}}}: Always-accurate
approximate-private clipping.  \phantomsection\label{alg:approx-clip-sum}

\uline{Input}: A database $D\in\{0,\ldots,U\}^n$; public parameters
$U,n,\varepsilon,\delta,\beta$ satisfying the hypotheses of
\Cref{thm:main}, with $\delta>0$.

\uline{Output}: A real estimate of $S(D)$.

\uline{Guarantee}: $(\varepsilon,\delta)$-privacy and deterministic
normalized error $O(B_{\varepsilon,\delta})$.

\begin{enumerate}[leftmargin=*,nosep]
  \item If $K_{\mathrm{app}}>n$, return $0$.
  \item Form
  $Y=\operatorname{Top}_{K_{\mathrm{app}}}
     (\{b(x_i):i\in[n]\})$ and compute
  $z=\text{\textsc{SureInterior}}(Y,\varepsilon/2,\delta/2)$.
  \item Put $t=2^z$ and return
  $\text{\textsc{TruncatedLaplaceRelease}}
  (D,Q_t;\Delta=t,\alpha=\varepsilon/2,\zeta=\delta/2)$.
\end{enumerate}
\end{algbox}

\begin{lemma}[Approximate-private clipping guarantee]
\label{lem:approx-upper}
Under the hypotheses of \Cref{lem:parameter-comparison},
\hyperref[alg:approx-clip-sum]{\textsc{ApproxClipSum}} is
$(\varepsilon,\delta)$-differentially private and satisfies, for every $D$,
\[
 |\text{\textsc{ApproxClipSum}}(D)-S(D)|
 \le\frac{C m(D)}\varepsilon B_{\varepsilon,\delta}
\]
on every execution.  It runs in time polynomial in
$n$, $\log U$, $\log(1/\delta)$, and $\log(1/\varepsilon)$ in the exact-real
model.
\end{lemma}

\begin{proof}
If $K_{\mathrm{app}}>n$, the output is data independent and
\[
 \frac{\varepsilon|S(D)|}{m(D)}
 \le\varepsilon n<\varepsilon K_{\mathrm{app}}
 =O(G)=O(B_{\varepsilon,\delta})
\]
by \Cref{lem:parameter-comparison}.

Suppose $K_{\mathrm{app}}\le n$.  By
\Cref{lem:order-truncation-stability,lem:sure-interior}, the selector is
$(\varepsilon/2,\delta/2)$-private.  Conditional on every fixed selected
threshold $t$, \Cref{lem:clipped-sensitivity,lem:truncated-laplace} make the
release $(\varepsilon/2,\delta/2)$-private.  Adaptive composition proves
overall $(\varepsilon,\delta)$-privacy.

The selected $z$ always lies between the minimum and maximum of $Y$.
Thus \Cref{lem:dyadic-clipping} bounds the clipping bias by
$K_{\mathrm{app}}m(D)$ and gives $t<2m(D)$.  The added noise has magnitude at
most
\[
 \begin{aligned}
 R\left(\frac\varepsilon2,\frac\delta2,t\right)
 &=\frac{2t}{\varepsilon}
   \log\left(1+\frac{e^{\varepsilon/2}-1}{\delta}\right)\\
 &\le\frac{4m(D)}\varepsilon B_{\varepsilon,\delta}.
 \end{aligned}
\]
Consequently,
\[
 \frac{\varepsilon
 |\text{\textsc{ApproxClipSum}}(D)-S(D)|}{m(D)}
 \le\varepsilon K_{\mathrm{app}}+4B_{\varepsilon,\delta}
 =O(G+B_{\varepsilon,\delta})
 =O(B_{\varepsilon,\delta}),
\]
deterministically.

For the resource bound, computing and sorting the scales costs
$O(n\log n)$.  The selector makes two calls to
\textsc{OneLevelPrefix}, and evaluating $Q_t$ costs $O(n)$; the bounded-noise
release adds $O(1)$ exact-real operations.  Here $K_{\mathrm{app}}\le n$,
$d=L_U$, and the base parameters are
\[
 \left(\alpha_0,\zeta_0,\eta\right)
 =\left(\frac\varepsilon8,\frac\delta{32},
        \frac{\min\{\varepsilon,\delta\}}{128}\right).
\]
The claimed polynomial follows from
\Cref{lem:one-level-prefix,lem:sure-interior}, since
$L_U=O(1+\log U)$ and
$\log(1/\eta)=O(1+\log(1/\varepsilon)+\log(1/\delta))$.
\end{proof}

\subsection{The matching lower bound}

For a fixed mechanism $M$, abbreviate its worst normalized error by
\[
 R(M)=\sup_{D\in\{0,\ldots,U\}^n}
 \frac{\varepsilon\operatorname{err}_{\beta}(M,D)}{m(D)}.
\]
The first packing places geometrically separated databases around the
all-zero database and retains the exact additive loss from group privacy.

\begin{lemma}[Geometric star-packing inequality]
\label{lem:star-packing}
Let $M$ be $(\varepsilon,\delta)$-differentially private, let $r>R(M)$, put
$a=r/\varepsilon$ and $k=\lfloor2a\rfloor+1$, and define
\[
 h=1+\lfloor\log_4U\rfloor,
 \qquad
 \gamma=
 \begin{cases}
  \delta/(e^\varepsilon-1),&\delta>0,\\
  0,&\delta=0.
 \end{cases}
\]
If $k\le n$, then
\begin{equation}
 e^{-k\varepsilon}
 \le\frac{\beta/h+\gamma}{1-\beta+\gamma},
 \qquad
 k\varepsilon
 \ge\log\frac{1-\beta+\gamma}{\beta/h+\gamma}.
 \label{eq:exact-star}
\end{equation}
\end{lemma}

\begin{proof}
For $j=0,\ldots,h-1$, let $v_j=4^j\le U$.  Let $D_\circ$ be the all-zero
database, and let $D_j$ contain $k$ copies of $v_j$ and $n-k$ zeroes.  Since
$r>R(M)$, the mechanism assigns probability at least $1-\beta$ to
\[
 I_\circ=[-a,a]
 \quad\text{under }D_\circ,
 \qquad
 I_j=[(k-a)v_j,(k+a)v_j]
 \quad\text{under }D_j.
\]
Indeed, each displayed radius is strictly larger than
$\operatorname{err}_{\beta}(M,D)$ for the corresponding database, and the
set of radii attaining probability at least $1-\beta$ is upward closed.
These intervals are pairwise disjoint.  The inequality $k>2a$ separates
$I_\circ$ from $I_0$, and consecutive geometric intervals are disjoint
because
\[
 (k+a)v_j<(k-a)v_{j+1}
 \quad\Longleftrightarrow\quad 5a<3k,
\]
which follows from $k>2a$.

The replacement distance from $D_\circ$ to $D_j$ is $k$.  Applying
\Cref{eq:group-privacy} from $D_j$ toward $D_\circ$ to the event $I_j$
and rearranging exactly yields
\[
 \Pr[M(D_\circ)\in I_j]
 \ge e^{-k\varepsilon}(1-\beta+\gamma)-\gamma.
\]
The interval $I_\circ$ already has probability at least $1-\beta$ under
$D_\circ$, so the other disjoint intervals have total probability at most
$\beta$.  Summing the displayed lower bound on
$\Pr[M(D_\circ)\in I_j]$ over $j$ proves both forms in
\Cref{eq:exact-star}.
\end{proof}

The logarithmic packing alone loses an additive constant.  A three-point
count packing supplies a uniform positive floor before those constants are
absorbed.

\begin{lemma}[Uniform constant error floor]
\label{lem:constant-floor}
If $\varepsilon n\ge64A_{U,\beta}$ and $\delta\le n^{-2}$, then every
$(\varepsilon,\delta)$-differentially private mechanism $M$ satisfies
$R(M)\ge1/20$.
\end{lemma}

\begin{proof}
Suppose that $R(M)<1/20$, and choose $R(M)<r<1/20$.  Put
\[
 a=\frac r\varepsilon,
 \qquad
 k=\lfloor2a\rfloor+1,
 \qquad
 \gamma=
 \begin{cases}
  \delta/(e^\varepsilon-1),&\delta>0,\\
  0,&\delta=0.
 \end{cases}
\]
Then
\[
 k<\frac{0.1}{\varepsilon}+1\le\frac{1.1}{\varepsilon}.
\]
Since $A_{U,\beta}\ge1+\log6$, the sample hypothesis gives
$n>2.2/\varepsilon>2k$.

For $q\in\{0,k,2k\}$, let $E_q$ contain $q$ ones and $n-q$ zeroes, and set
$J_q=[q-a,q+a]$.  These intervals are disjoint because $k>2a$.  Moreover,
$m(E_q)=1$ and $a>\operatorname{err}_{\beta}(M,E_q)$; by upward closure of
the admissible quantile radii, $J_q$ has probability at least $1-\beta$
under $E_q$.  Each of $E_0$ and $E_{2k}$ is at replacement distance $k$ from
$E_k$.  Group privacy in \Cref{eq:group-privacy}, rearranged toward $E_k$,
gives
\[
 \Pr[M(E_k)\in J_0],\ \Pr[M(E_k)\in J_{2k}]
 \ge e^{-k\varepsilon}(1-\beta)-\gamma.
\]
Since $J_k$ itself has mass at least $1-\beta$ under $E_k$, disjointness
implies
\begin{equation}
 \beta\ge2e^{-k\varepsilon}(1-\beta)-2\gamma.
 \label{eq:floor-contradiction}
\end{equation}
Here $k\varepsilon<1.1$.  Also, using
$e^\varepsilon-1\ge\varepsilon$, $\delta\le n^{-2}$, and the sample
hypothesis,
\[
 \gamma\le\frac1{\varepsilon n^2}
 \le\frac1{(\varepsilon n)^2}<\frac1{18000}.
\]
Because $1-\beta\ge2/3$, the right-hand side of
\Cref{eq:floor-contradiction} is strictly greater than
$(4/3)e^{-1.1}-2/18000>1/3\ge\beta$, a contradiction.
\end{proof}

Combining the exact star inequality with the constant floor yields the full
rate.  The case in which the packing would require too many records is
handled directly.

\begin{lemma}[Matching minimax lower bound]
\label{lem:full-lower}
If $\varepsilon n\ge64A_{U,\beta}$ and $\delta\le n^{-2}$, then
\[
 \rho^*_{\varepsilon,\delta,\beta}(n,U)
 \ge\frac1{200}\min\{A_{U,\beta},B_{\varepsilon,\delta}\}.
\]
\end{lemma}

\begin{proof}
Fix an $(\varepsilon,\delta)$-differentially private mechanism $M$.  If
$R(M)=+\infty$, there is nothing to prove.  Otherwise choose $r>R(M)$ and
put
\[
 a=\frac r\varepsilon,
 \qquad
 k=\lfloor2a\rfloor+1,
 \qquad
 T=\min\{A_{U,\beta},B_{\varepsilon,\delta}\},
 \qquad
 h=1+\lfloor\log_4U\rfloor.
\]
If $k>n$, then $\lfloor2a\rfloor\ge n$, and hence
\[
 r\ge\frac{\varepsilon n}{2}
 \ge32A_{U,\beta}\ge32T.
\]
Assume now that $k\le n$.  From \Cref{lem:star-packing}, weakening the
numerator to $1-\beta$, using $1-\beta\ge2/3$, and applying
$x+y\le2\max\{x,y\}$ gives
\begin{equation}
 k\varepsilon
 \ge\min\left\{
  \log\frac h\beta,
  \log\frac{e^\varepsilon-1}{\delta}
 \right\}-\log3,
 \label{eq:two-log-lower}
\end{equation}
where the second logarithm is $+\infty$ when $\delta=0$.

Writing $q=\lfloor(\log_2U)/2\rfloor$, we have $h=q+1$ and
$L_U\le2q+3\le3h$.  Therefore
\[
 \log(h/\beta)\ge A_{U,\beta}-1-\log3.
\]
When $\delta>0$, the standing hypotheses imply
\[
 \frac{e^\varepsilon-1}{\delta}
 \ge\varepsilon n^2
 =\frac{(\varepsilon n)^2}{\varepsilon}>1,
\]
and consequently
\[
 \log\frac{e^\varepsilon-1}{\delta}
 \ge B_{\varepsilon,\delta}-1-\log2.
\]
Substitution in \Cref{eq:two-log-lower} yields
$k\varepsilon\ge T-1-2\log3$.  On the other hand,
$k\varepsilon\le2r+\varepsilon\le2r+1$, so
\begin{equation}
 r\ge\frac{T-2-2\log3}{2}.
 \label{eq:lower-up-to-constant}
\end{equation}

By \Cref{lem:constant-floor}, $R(M)\ge1/20$.  If
$T\ge2(2+2\log3)$, \Cref{eq:lower-up-to-constant} gives $r\ge T/4$.
Otherwise $T<8.4$ and $R(M)\ge1/20>T/200$.  Letting
$r\downarrow R(M)$ shows $R(M)\ge T/200$.  Taking the infimum over all
private mechanisms proves the lemma.
\end{proof}

\subsection{Public branch selection and conclusion}

The final mechanism chooses between the two upper constructions using only
public parameters.  Thus the branch decision consumes no privacy budget.

\begin{algbox}
\textbf{\uline{\textsc{BranchClipSum}}}: Instance-optimal bounded
summation.  \phantomsection\label{alg:branch-clip-sum}

\uline{Input}: A database $D\in\{0,\ldots,U\}^n$; public parameters
$U,n,\varepsilon,\delta,\beta$ satisfying the hypotheses of
\Cref{thm:main}.

\uline{Output}: A real estimate of $S(D)$.

\uline{Guarantee}: $(\varepsilon,\delta)$-privacy and normalized
$\beta$-error $O(\min\{A_{U,\beta},B_{\varepsilon,\delta}\})$.

\begin{enumerate}[leftmargin=*,nosep]
  \item Compute the public quantities $A_{U,\beta}$ and
  $B_{\varepsilon,\delta}$, using $B_{\varepsilon,0}=+\infty$.
  \item If $\delta=0$ or $A_{U,\beta}\le B_{\varepsilon,\delta}$, return
  $\text{\textsc{PureClipSum}}(D)$.
  \item Otherwise return $\text{\textsc{ApproxClipSum}}(D)$.
\end{enumerate}
\end{algbox}

\begin{lemma}[Best-branch privacy, accuracy, and resources]
\label{lem:branch-wrapper}
Under the hypotheses of \Cref{thm:main},
\hyperref[alg:branch-clip-sum]{\textsc{BranchClipSum}} is
$(\varepsilon,\delta)$-differentially private and, for every $D$,
\[
 \Pr\left[
  |\text{\textsc{BranchClipSum}}(D)-S(D)|
  \le\frac{C m(D)}\varepsilon
       \min\{A_{U,\beta},B_{\varepsilon,\delta}\}
 \right]\ge1-\beta.
\]
On the pure branch its running time is $O(n\log n+L_U)$; on the approximate
branch it is polynomial in $n$, $\log U$, $\log(1/\delta)$, and
$\log(1/\varepsilon)$.
\end{lemma}

\begin{proof}
The branch decision depends only on public parameters.  On the pure branch,
\Cref{lem:pure-upper} gives $(\varepsilon,0)$-privacy and normalized
$\beta$-error $O(A_{U,\beta})$, and here
$A_{U,\beta}=\min\{A_{U,\beta},B_{\varepsilon,\delta}\}$.  On the
approximate branch, \Cref{lem:approx-upper} gives
$(\varepsilon,\delta)$-privacy and deterministic normalized error
$O(B_{\varepsilon,\delta})$, and here
$B_{\varepsilon,\delta}<A_{U,\beta}$.  Exactly one branch is invoked, so its
resource bound is inherited unchanged.
\end{proof}

By \Cref{lem:parameter-bookkeeping}, the sample-size condition in
\Cref{thm:main} implies the premise of \Cref{lem:full-lower}, which proves
the left inequality in \Cref{eq:main-rate}.  The right inequality and the
claimed output and resource guarantees follow from
\Cref{lem:branch-wrapper}.  This completes the proof of \Cref{thm:main}.

\bibliographystyle{alpha}
\bibliography{references}

\end{document}
